\documentclass[UTF8,a4paper,12pt]{ctexart}

\usepackage{geometry}
\usepackage{graphicx}
\usepackage{booktabs}
\usepackage{longtable}
\usepackage{array}
\usepackage{xcolor}
\usepackage{hyperref}
\usepackage{enumitem}
\usepackage{listings}
\usepackage{float}
\usepackage{caption}

\geometry{left=2.6cm,right=2.6cm,top=2.5cm,bottom=2.5cm}
\hypersetup{
  colorlinks=true,
  linkcolor=blue,
  urlcolor=blue,
  citecolor=blue
}

\setlist[itemize]{leftmargin=2em}
\setlist[enumerate]{leftmargin=2em}
\renewcommand{\arraystretch}{1.25}
\graphicspath{{../../test-results/}}

\lstdefinestyle{codestyle}{
  basicstyle=\ttfamily\small,
  breaklines=true,
  frame=single,
  columns=fullflexible,
  keepspaces=true,
  showstringspaces=false,
  backgroundcolor=\color{gray!6},
  rulecolor=\color{gray!45}
}
\lstset{style=codestyle}

\title{\textbf{基于 Python Playwright 的百度地图 Web 自动化测试工程实验报告}}
\author{Baidu Map Pytest Playwright Test Lab}
\date{\today}

\begin{document}

\maketitle
\thispagestyle{empty}

\begin{abstract}
本实验围绕百度地图 Web 端构建了一套基于 Python Playwright、pytest、
pytest-playwright 与 uv 的端到端自动化测试工程。实验覆盖首页加载、搜索流程、
路线规划、定位权限、移动端视口、HTML 报告、截图、Trace 以及受控失败证据。
由于百度地图属于公开复杂 Web 地图应用，实验特别关注动态 DOM、异步加载、地图
瓦片渲染、反自动化安全验证、外部网络波动等真实约束。最终稳定套件包含 8 个
测试用例，连续三轮重复执行均通过，最新记录为 8 passed in 37.68s。
\end{abstract}

\newpage
\tableofcontents
\newpage

\section{概述}

\subsection{实验背景}

百度地图 Web 端是典型的复杂前端应用。与普通表单页面相比，地图类应用具有更
多动态行为：地图瓦片按视口加载，搜索和路线结果依赖网络响应，页面包含广告、
定位权限、复杂输入联想、地图拖拽缩放、路线方案切换等交互。此类系统适合作为
现代 Web UI 自动化测试工具的实验对象。

本实验选择 Playwright for Python 作为核心自动化工具，pytest 作为测试组织与
执行框架，pytest-playwright 作为浏览器 fixture 与命令行参数支持，uv 作为
Python 环境和依赖管理工具。实验目标不是对百度地图进行数据采集，也不是绕过
公开站点安全验证，而是在真实公开网站约束下验证自动化测试工程的可执行性、
可诊断性和可维护性。

\subsection{实验目标}

本实验的工程目标包括：

\begin{itemize}
  \item 搭建可复现的 Python Playwright 自动化测试工程。
  \item 使用 Page Object Model 封装百度地图页面操作。
  \item 使用 pytest 标记和 fixture 组织首页、搜索、路线、定位、移动端测试。
  \item 使用 HTML 报告、截图和 Trace 形成可检查证据链。
  \item 在不绕过安全验证的前提下记录公开站点反自动化约束。
  \item 分析 Playwright 在复杂地图类 Web 应用测试中的优势与边界。
\end{itemize}

\subsection{测试对象与实验边界}

被测对象为百度地图 Web 端，地址为：

\begin{center}
\url{https://map.baidu.com}
\end{center}

实验明确排除以下内容：

\begin{itemize}
  \item 不进行高频请求、批量采集、压力测试或安全扫描。
  \item 不绕过滑块、图像旋转、扫码等百度安全验证。
  \item 不验证地图数据的绝对准确性。
  \item 不断言实时路线分钟数、公里数、红绿灯数量、价格、路况或路线几何。
  \item 不测试百度地图移动 App，仅测试 Web 页面在移动视口下的基本可用性。
\end{itemize}

\section{测试工具环境建立}

\subsection{工程环境}

项目采用 uv 管理 Python 环境和依赖，当前工程配置位于 \path{pyproject.toml}。
核心环境如下：

\begin{table}[H]
\centering
\caption{实验工具链}
\begin{tabular}{ll}
\toprule
工具 & 用途 \\
\midrule
Python 3.12 & 测试代码运行环境 \\
uv & 虚拟环境、依赖和命令运行管理 \\
Playwright for Python & 浏览器自动化核心能力 \\
pytest & 测试用例组织、断言和执行 \\
pytest-playwright & Playwright 浏览器 fixture 与 CLI 集成 \\
pytest-html & 生成 HTML 测试报告 \\
ruff & Python 静态检查与格式检查 \\
pre-commit & 提交前基础质量检查 \\
just & 项目常用命令封装 \\
\bottomrule
\end{tabular}
\end{table}

\subsection{依赖和浏览器安装}

项目初始化和浏览器安装命令如下：

\begin{lstlisting}[language=bash]
uv sync
uv run playwright install
\end{lstlisting}

由于 Codex sandbox 中 uv 默认 cache 可能不可写，项目提供了
\path{scripts/agent-env.sh}，通过 \texttt{UV\_CACHE\_DIR=.cache/uv} 将 cache
重定向到工作区内。agent-safe 命令由 \path{justfile} 封装：

\begin{lstlisting}[language=bash]
just agent-ruff
just agent-test
just agent-failure-demo
\end{lstlisting}

\subsection{pytest 与 Playwright 配置}

pytest 默认配置包括：

\begin{itemize}
  \item 测试目录：\path{tests}
  \item 默认浏览器：Chromium
  \item 失败时保留 Trace、截图和视频
  \item Playwright 自动产物目录：\path{test-results/playwright}
  \item HTML 报告路径：\path{reports/report.html}
\end{itemize}

项目使用如下 pytest mark 区分用例类型：

\begin{table}[H]
\centering
\caption{pytest mark 设计}
\begin{tabular}{ll}
\toprule
标记 & 含义 \\
\midrule
\texttt{smoke} & 核心冒烟检查 \\
\texttt{search} & 百度地图搜索流程 \\
\texttt{route} & 路线规划流程 \\
\texttt{geo} & 定位权限和模拟坐标 \\
\texttt{mobile} & 移动端视口测试 \\
\texttt{slow} & 可能受网络或地图加载影响较大的测试 \\
\bottomrule
\end{tabular}
\end{table}

\section{测试工具的功能和使用流程}

\subsection{Playwright 页面操作与 Locator}

Playwright 通过 \texttt{Page} 对象驱动浏览器页面。项目中将直接页面操作封装在
\texttt{BaiduMapPage} 中，测试用例只表达业务流程。例如首页测试只需要执行：

\begin{lstlisting}[language=Python]
baidu_map.goto()
baidu_map.expect_page_loaded()
baidu_map.screenshot("homepage")
\end{lstlisting}

Locator 选择优先使用 role、文本、placeholder 等语义信息。例如搜索输入框使用：

\begin{lstlisting}[language=Python]
page.get_by_role("textbox", name="搜地点、查公交、找路线")
\end{lstlisting}

路线入口和路线输入框则根据人工 codegen 观察确定：

\begin{lstlisting}[language=Python]
page.locator(".searchbox-content-button")
page.get_by_role("textbox", name="输入起点或在图区上选点")
page.get_by_role("textbox", name="输入终点")
\end{lstlisting}

\subsection{自动等待与弹性断言}

Playwright 的 \texttt{expect(locator).to\_be\_visible()} 会自动等待元素达到目标
状态。对公开地图应用而言，断言不能过于精确。本实验采用如下策略：

\begin{itemize}
  \item 首页只断言 body 可见。
  \item 搜索流程不强制断言最终结果页必定出现，而是分类记录结果、安全验证、
  URL 跳转、输入被清空等可观察状态。
  \item 路线规划只断言路线方案区域和路线细节区域出现，不断言实时数值。
  \item geolocation 与 mobile 只验证页面稳定可打开。
\end{itemize}

\subsection{截图、Trace 与 HTML 报告}

项目通过 \texttt{BaiduMapPage.screenshot()} 将关键证据写入
\path{test-results/}。pytest-playwright 负责在失败时保留 Trace、截图和视频；
pytest-html 负责生成 \path{reports/report.html}。

受控失败证据由以下命令生成：

\begin{lstlisting}[language=bash]
just agent-failure-demo
\end{lstlisting}

Trace 查看方式如下：

\begin{lstlisting}[language=bash]
uv run playwright show-trace \
  test-results/failure-demo/failure-demo-trace.zip
\end{lstlisting}

\subsection{Browser Context、定位权限与移动端}

Playwright Browser Context 允许每个测试使用独立的权限、定位、语言、时区和视口。
本实验使用 context 模拟北京定位：

\begin{lstlisting}[language=Python]
browser.new_context(
    geolocation=GEO_LOCATIONS["beijing"],
    permissions=["geolocation"],
    locale="zh-CN",
    timezone_id="Asia/Shanghai",
)
\end{lstlisting}

移动端 smoke 使用 393 x 852 viewport，并开启 \texttt{is\_mobile} 和
\texttt{has\_touch}，用于验证百度地图 Web 页面在移动视口下可以稳定打开。

\subsection{pytest 作为辅助测试工具}

pytest 负责测试发现、用例组织、mark 管理和报告输出。项目按照功能拆分为：

\begin{itemize}
  \item \path{tests/test_home.py}
  \item \path{tests/test_search.py}
  \item \path{tests/test_route.py}
  \item \path{tests/test_geolocation.py}
  \item \path{tests/test_mobile.py}
\end{itemize}

\section{被测软件介绍}

百度地图 Web 端提供搜索、路线、定位、地图浏览和图层切换等能力。本实验重点关注
以下模块：

\begin{longtable}{p{0.22\linewidth}p{0.68\linewidth}}
\caption{被测功能模块} \\
\toprule
模块 & 特征 \\
\midrule
\endfirsthead
\toprule
模块 & 特征 \\
\midrule
\endhead
首页加载 & 打开地图首页，加载搜索框、地图区域和基础控件。 \\
搜索 & 输入地点关键词，触发联想、搜索提交、安全验证或地图跳转。 \\
路线规划 & 打开路线面板，选择驾车模式，输入起点和终点并生成路线方案。 \\
定位 & 浏览器请求或模拟 geolocation 权限，页面需稳定处理定位上下文。 \\
移动端视口 & 在移动 viewport 和 touch context 下验证页面可打开。 \\
地图交互 & 地图瓦片、缩放、拖拽、右侧广告卡片和动态 UI 共同影响自动化稳定性。 \\
\bottomrule
\end{longtable}

百度地图页面存在明显的反自动化与安全验证机制。人工 codegen 观察显示，提交搜索
可能触发“百度安全验证”和“请完成下方验证后继续操作”，验证方式包括图像旋转、
滑块和扫码验证。实验不绕过这些机制，而是将其作为测试观察结果。

\section{测试计划与测试用例设计}

\subsection{测试范围}

本实验最终实现了 8 个稳定用例，覆盖首页、搜索、路线、定位和移动端。

\begin{table}[H]
\centering
\caption{稳定测试套件组成}
\begin{tabular}{lll}
\toprule
模块 & 用例数 & 说明 \\
\midrule
首页 smoke & 1 & 首页可打开并截图 \\
搜索流程 & 3 & 输入、联想状态分类、结果状态分类 \\
路线规划 & 2 & 路线面板 smoke、驾车路线 baseline \\
定位权限 & 1 & 北京 geolocation context smoke \\
移动端 & 1 & 移动 viewport smoke \\
\midrule
合计 & 8 & 最新稳定运行全部通过 \\
\bottomrule
\end{tabular}
\end{table}

\subsection{测试数据}

测试数据集中维护在 \path{data/test_data.py}：

\begin{itemize}
  \item 搜索关键词：天安门广场、北京大学、清华大学、北京南站等。
  \item 异常输入：空字符串、特殊字符、不存在地点字符串。
  \item 路线用例：北京南站到天安门广场、北京大学到清华大学等。
  \item 定位坐标：北京、上海、广州的公开城市中心坐标。
\end{itemize}

最终稳定路线 baseline 使用：

\begin{center}
北京南站 $\rightarrow$ 天安门广场
\end{center}

\subsection{用例设计}

\begin{longtable}{p{0.16\linewidth}p{0.16\linewidth}p{0.58\linewidth}}
\caption{最终自动化用例} \\
\toprule
用例编号 & mark & 验证目标 \\
\midrule
\endfirsthead
\toprule
用例编号 & mark & 验证目标 \\
\midrule
\endhead
TC-HOME-001 & smoke & 首页 body 可见并生成截图，对应
\path{tests/test_home.py}。 \\
TC-SEARCH-001 & search & 搜索框接受“北京大学”输入。 \\
TC-SEARCH-002 & search & 对联想、输入保留、输入清空、安全验证等状态分类。 \\
TC-SEARCH-003 & search & 对搜索提交和已知结果 URL 的结果状态分类。 \\
TC-ROUTE-001 & route & 路线入口可打开路线面板。 \\
TC-ROUTE-002 & route & 驾车路线 baseline 出现方案和步骤详情。 \\
TC-GEO-001 & geo & 北京模拟定位上下文可以打开地图。 \\
TC-MOBILE-001 & mobile & 移动端视口可以打开地图。 \\
\bottomrule
\end{longtable}

完整 pytest 函数名保留在 \path{tests/} 目录中，报告正文使用用例编号提高表格
可读性。

\subsection{反自动化约束下的分类断言}

搜索流程不把“出现百度安全验证”视为自动化失败，因为这是公开站点的正常保护
机制。搜索测试的分类结果包括：

\begin{itemize}
  \item \texttt{suggestion}：出现搜索联想。
  \item \texttt{security\_challenge}：出现百度安全验证。
  \item \texttt{map\_navigation}：headless 下提交后发生地图 URL 跳转。
  \item \texttt{input\_cleared}：输入被页面清空。
  \item \texttt{input\_retained\_without\_suggestion}：输入保留但无联想。
  \item \texttt{partial\_input\_retained}：只保留部分输入。
\end{itemize}

\section{测试工具的应用}

\subsection{项目结构}

项目核心目录如下：

\begin{lstlisting}[language=bash]
data/              # 测试数据
docs/              # 设计书、最终报告、locator 观察
pages/             # Page Object
reports/           # pytest-html 报告和 LaTeX 报告工程
scripts/           # agent-safe uv wrapper、失败证据脚本
test-results/      # 截图、trace、失败证据
tests/             # pytest 用例
\end{lstlisting}

\subsection{核心代码}

页面操作集中在 \texttt{BaiduMapPage}。例如路线规划 baseline 使用以下 Page
Object 方法：

\begin{lstlisting}[language=Python]
def plan_drive_route(self, start: str, end: str) -> None:
    self.open_route_panel()
    self.drive_mode_tab.click()
    self.route_start_input.click()
    self.route_start_input.press_sequentially(start, delay=120)
    self.route_start_input.press("Enter")
    route_end_input = self.route_end_input.first
    route_end_input.click()
    route_end_input.press_sequentially(end, delay=120)
    route_end_input.press("Enter")
\end{lstlisting}

稳定性修复中，页面导航增加了最多 3 次重试，用于缓解公开站点偶发
\texttt{net::ERR\_NETWORK\_CHANGED}：

\begin{lstlisting}[language=Python]
def goto_url(self, url: str) -> None:
    last_error = None
    for _ in range(3):
        try:
            self.page.goto(url, wait_until="domcontentloaded")
            return
        except PlaywrightError as error:
            last_error = error
            self.page.wait_for_timeout(1000)
    raise last_error
\end{lstlisting}

\subsection{运行命令}

常用执行命令：

\begin{lstlisting}[language=bash]
just agent-ruff
just agent-test
just agent-failure-demo
\end{lstlisting}

报告编译命令：

\begin{lstlisting}[language=bash]
cd reports/latex
make pdf
\end{lstlisting}

\subsection{执行结果}

最终稳定套件最新结果：

\begin{lstlisting}
8 passed in 37.68s
\end{lstlisting}

连续三轮重复执行结果：

\begin{table}[H]
\centering
\caption{重复执行稳定性}
\begin{tabular}{lll}
\toprule
轮次 & 结果 & 耗时 \\
\midrule
1 & 8 passed & 39.94s \\
2 & 8 passed & 41.30s \\
3 & 8 passed & 37.68s \\
\bottomrule
\end{tabular}
\end{table}

重复执行前曾发现两个稳定性问题：一次网络变化导致导航失败，一次路线结果断言被
地图比例尺中的“公里”误触发。修复后连续三次通过，说明当前稳定套件具备基本
可重复执行能力。

\subsection{截图证据}

\begin{figure}[H]
\centering
\includegraphics[width=0.9\linewidth]{homepage.png}
\caption{首页 smoke 截图}
\end{figure}

\begin{figure}[H]
\centering
\includegraphics[width=0.9\linewidth]{search-input-keyword.png}
\caption{搜索输入截图}
\end{figure}

\begin{figure}[H]
\centering
\includegraphics[width=0.9\linewidth]{search-known-result-page-security_challenge.png}
\caption{已知搜索结果 URL 触发百度安全验证}
\end{figure}

\begin{figure}[H]
\centering
\includegraphics[width=0.9\linewidth]{route-result-beijing_south_to_tiananmen.png}
\caption{驾车路线方案区域}
\end{figure}

\begin{figure}[H]
\centering
\includegraphics[width=0.9\linewidth]{route-detail-beijing_south_to_tiananmen.png}
\caption{驾车路线细节步骤}
\end{figure}

\begin{figure}[H]
\centering
\includegraphics[width=0.9\linewidth]{geolocation-beijing.png}
\caption{北京模拟定位上下文}
\end{figure}

\begin{figure}[H]
\centering
\includegraphics[width=0.5\linewidth]{mobile-homepage.png}
\caption{移动端视口 smoke}
\end{figure}

\subsection{HTML 报告与失败 Trace}

HTML 报告路径为：

\begin{center}
\path{reports/report.html}
\end{center}

受控失败证据路径为：

\begin{itemize}
  \item \path{test-results/failure-demo/failure-demo-missing-text.png}
  \item \path{test-results/failure-demo/failure-demo-trace.zip}
  \item \path{test-results/failure-demo/failure-summary.txt}
\end{itemize}

该失败演示故意等待不存在的文本，用来展示 Playwright 在失败时提供的截图和
Trace 诊断能力。该失败不属于稳定 pytest 套件，不影响 8 个稳定用例的通过率。

\section{总结}

\subsection{Playwright、pytest 与 uv 的优点}

\begin{itemize}
  \item Playwright 的 Locator、自动等待、Browser Context、截图和 Trace 能力
  适合复杂 Web UI 的端到端验证。
  \item pytest 提供了清晰的用例组织、mark 分类、fixture 复用和 HTML 报告扩展。
  \item uv 提供了快速、可复现的 Python 依赖和命令运行环境。
  \item Page Object Model 降低了测试用例与页面细节的耦合。
\end{itemize}

\subsection{不足与经验}

\begin{itemize}
  \item 公开地图站点存在安全验证，完整搜索流程不能强行自动化。
  \item Headed 和 headless 行为存在差异，例如搜索联想和路线候选项表现不同。
  \item 外部网络波动会影响端到端测试稳定性，需要重试和失败证据。
  \item 地图瓦片、实时路况、路线时长等动态数据不适合精确断言。
\end{itemize}

\subsection{与 JMeter、Wireshark、Postman 的互补关系}

Playwright 关注用户视角的浏览器端流程，适合验证页面交互和端到端用户路径。
JMeter 更适合接口或服务端压力测试；Wireshark 更适合网络抓包和协议层分析；
Postman 更适合 API 调试、接口集合管理和接口回归测试。对于复杂 Web 应用，四类
工具并非替代关系，而是覆盖不同测试层次：Playwright 验证真实浏览器行为，
Postman 验证接口契约，JMeter 验证负载能力，Wireshark 分析网络通信细节。

\subsection{最终结论}

本实验完成了一套可执行、可诊断、可复现的百度地图 Web 自动化测试工程。最终
稳定套件覆盖首页、搜索、路线、定位和移动端，形成 HTML 报告、截图和 Trace
证据。实验表明，Python Playwright + pytest + uv 适合复杂地图类 Web 应用的
工程化自动化测试，但测试设计必须尊重公开站点安全边界，避免高频请求和过度
精确断言。

\end{document}
